\scanpage{446}
\mspage{358}
\uline{Définition} 4.13.35. \unsure{Soient|Soit} $M$ une catégorie accessible, stable par petites $\underrightarrow{\lim}$ filtrantes, \deleted{Une partie} $S$ une partie de $\mathrm{Ob}\,M$, $\underline S$ la sous-catégorie \deleted{\unsure{str.}} pleine correspondante de $M$. On dit que $S$ est $L$-\uline{accessible}, \deleted{\unsure{si}} que $\underline S$ est $L$-accessible dans $M$, si $\underline S$ est accessible dans $M$, et de plus stable par \inserted{petites} $\underrightarrow{\lim}$ filtrantes.

\margin{\uline{NB} \deleted{Les} \inserted{L'ens des} parties \uline{$L$}-accessibles de $\mathrm{Ob}\,M$ est stable par petites intersections, et par réunions \uline{finies}. L'image inverse d'une partie \unsure{$L$|$L_0$}-accessible par un $L$-foncteur\textsuperscript{*} $f:M\unsure{\text{---}|\longrightarrow}N$ de cat. acc\,; \deleted{\unsure{et}} est $L$-accessible.}
% slanted
\margin{un \uline{$L_0$-foncteur} : qui commute aux petites $\underrightarrow{\lim}$ filtrantes.}

\uline{Proposition} 4.13.36. Soient $M,S,\underline S$ comme ds 4.13.35, et soit $\pi$ un cardinal utile pour $M$. \deleted{Conditions} \inserted{Supposons $\underline S$ \deleted{\unsure{$\pi$}} stable \deleted{par} dans $M$ par petites $\underrightarrow{\lim}$ filtrantes, ou} \deleted{équivalentes} seulement $\underline S(\pi)$ \inserted{$\overset{\mathrm{def}}{=}\underline S\cap M_\pi$} stable par petites $\underrightarrow{\lim}$ filtrantes \inserted{(\unsure{\deleted{dans}|dans} $M_\pi$)} de cardinal $\leq\pi$. Conditions équivalentes :

a) $\underline S$ est $\pi$-accessible, \deleted{et stable par petites $\underrightarrow{\lim}$ filtrantes} \deleted{\unsure{est est}} (\unsure{Il sera automatiquement|Il suivra toujours} stable par petites $\underrightarrow{\lim}$ filtrantes, i.e. $L$-accessible).

\deleted{a$'$) $\underline S$ est $\pi$-accessible, et $\underline S(\pi)\overset{\mathrm{def}}{=}\underline S\cap M_\pi$ est stable par $\underrightarrow{\lim}$ filtrantes de card $\leq\pi$, (propriété $L(\pi)$).}

b) \deleted{$\underline S(\pi)$ satisfait $L(\pi)$, et} Pour tt $x\in\mathrm{Ob}\,M$, on a $x\in\mathrm{Ob}\,\underline S$ sss $\underline S(\pi)_{/x}$ est cofinale ds $M_{\pi/x}$\unsure{,|.}

c) \deleted{Pour \unsure{tt} $\underline S(\pi)$ satisfait\,;} Pour tt $x\in\mathrm{Ob}\,M$, et tte représentation $x\simeq\underset{i\in I}{\underrightarrow{\lim}}\,x_i$
% archive sheet number: 445
